\chapter{Introduction}

\section{Context and Motivation}

Database management systems occupy a uniquely sensitive position in the modern software stack: they are responsible for storing, indexing, and serving the data on which virtually every application depends, from banking infrastructure to personal health records to critical national systems. A vulnerability in a database engine therefore does not merely threaten one application but potentially every application built on top of it. SQLite in particular stands out as the world's most widely deployed database engine, embedded in every major smartphone operating system, every mainstream web browser, the Python standard library, and countless embedded devices \cite{sqlite}. Its ubiquity makes it an especially consequential attack surface. Despite this importance, and despite extensive manual testing and code review, vulnerabilities continue to be discovered in SQLite at a steady pace. Between 2019 and 2020 alone, six Common Vulnerabilities and Exposures (CVEs) were filed against SQLite, affecting released versions 3.30.1 through 3.32.2: an integer overflow in the \texttt{INTERSECT} operator (CVE-2019-19646 \cite{cve201919646}), a stack overflow in expression parsing (CVE-2020-13434 \cite{cve202013434}), a null pointer dereference in the \texttt{isAggregateFunction} routine (CVE-2020-9327 \cite{cve20209327}), a null pointer dereference in the recursive \texttt{isAggregateFunction} (CVE-2020-13435 \cite{cve202013435}), a use-after-free triggered by the \texttt{ALTER TABLE} rename mechanism (CVE-2020-13871 \cite{cve202013871}), and a heap buffer overflow in multi-threaded \texttt{SELECT} with window functions (CVE-2020-15358 \cite{cve202015358}). The persistence of such defects in a heavily scrutinised, mature codebase underscores the need for more systematic automated approaches to vulnerability discovery.

Fuzzing has emerged as one of the most productive techniques for automated vulnerability detection, responsible for discovering thousands of bugs across the software industry \cite{fuzzingsurvey}. The central idea is to generate large volumes of semi-random inputs, feed them to a program under test, and observe whether the program crashes, hangs, or violates a safety property enforced by a dynamic analysis tool such as AddressSanitizer \cite{asan}. Coverage-guided greybox fuzzers like AFL \cite{afl} and libFuzzer \cite{libfuzzer} refine this by instrumenting the target binary to track which code branches are exercised by each input, then retaining inputs that increase coverage for further mutation. This feedback loop allows the fuzzer to progressively steer execution toward unexplored regions of the program. In practice, coverage-guided greybox fuzzing has proven highly effective for file parsers, network protocol implementations, and similar targets where input structure is loose or where arbitrary byte sequences can still reach meaningful program states. However, such byte-level mutation fuzzers encounter a fundamental obstacle when applied to programs that expect structured inputs, such as database engines that parse SQL. Random byte flips, bit-flips, and splice operations overwhelmingly produce inputs that are rejected at the parser stage, long before they can reach the deeper, semantically complex logic where vulnerabilities typically reside. The result is that the fuzzer saturates shallow code coverage quickly and fails to exercise the query optimiser, the window function engine, the virtual table subsystem, or other high-complexity components.

Grammar-based fuzzing addresses this problem by generating inputs from a formal specification of the target language, most commonly a context-free grammar, thereby guaranteeing syntactic validity and enabling the fuzzer to reach deeper program logic from the very first input. Nautilus \cite{nautilus} represents a particularly influential instance of this paradigm: it combines grammar-based generation with AFL-style coverage feedback, maintaining a corpus of grammar derivation trees rather than raw byte strings, and mutating those trees through subtree splicing, rule replacement, and random regeneration. This approach achieves substantially deeper code coverage on structured-input targets than byte-level fuzzers. Other grammar-aware fuzzers occupy adjacent points in the design space. Superion \cite{superion} extends AFL with grammar-aware mutations for JavaScript and XML targets. Grimoire \cite{grimoire} synthesises approximate structural patterns without requiring an explicit grammar specification, instead generalising from inputs that achieve new coverage. For relational databases specifically, Squirrel \cite{squirrel} generates SQL queries with explicit awareness of language validity constraints and coverage feedback, while SQLsmith \cite{sqlsmith} takes a simpler approach of generating random but syntactically valid SQL queries without coverage guidance.

A limitation shared by all of these grammar-based approaches, however, is the strategy used to select which production rule to apply at each step of input generation. Existing systems uniformly sample from all applicable rules with equal probability, meaning that a rule producing a common \texttt{SELECT} statement with a simple \texttt{WHERE} clause competes on equal footing with a rule producing a window function frame specification or a subquery inside a \texttt{CHECK} constraint. In a grammar with hundreds of rules, uniform sampling is necessarily diffuse: a large fraction of the mutation budget is consumed by rules that generate inputs which exercise already-covered code paths, yielding no new coverage signal and contributing nothing toward the discovery of new vulnerability classes. The uniform policy is stateless and memoryless, incapable of learning from the history of which rules have been productive. This inefficiency motivates the investigation of adaptive rule selection strategies that direct the mutation budget toward grammar rules with demonstrated potential to expand coverage.

\section{Problem Statement}

This thesis investigates whether adaptive rule selection can improve the effectiveness of grammar-based greybox fuzzing for vulnerability discovery in database management systems. The central research question is: can a multi-armed bandit algorithm \cite{exp3} that learns from coverage feedback to prioritise grammar production rules discover more diverse crash classes than uniform random sampling, when applied to a grammar-based fuzzer targeting SQLite?

The bandit hypothesis is that adaptive sampling, by concentrating mutation effort on rules that have historically expanded coverage, should guide the fuzzer toward unexplored program regions more efficiently than uniform sampling, thereby increasing both the speed and diversity of vulnerability discovery. The counter-hypothesis, which the experimental results of this thesis will ultimately support, is that such exploitation of historically productive rules may simultaneously reduce the structural diversity of the generated input corpus. Vulnerability discovery in mature software often requires rare, compositionally unusual inputs that exercise corner cases involving uncommon SQL constructs — precisely the inputs that an exploitation-focused bandit policy may systematically underweight. The interplay between coverage exploitation and structural diversity therefore constitutes the core tension this thesis examines.

\section{Objectives and Scope}

This thesis pursues three interconnected objectives, each building on the previous. The first objective is to design and implement a grammar-based greybox fuzzing system for SQLite that encodes SQL patterns known to trigger specific vulnerability classes without reproducing proof-of-concept exploit inputs verbatim. This requires a structural primitives methodology in which the grammar decomposes SQL into layered non-terminal categories — schema setup operations, stress queries involving complex joins and aggregations, validation operations that exercise integrity checking, and boundary function calls that probe edge cases in built-in SQL functions — such that the compositional combination of these primitives can independently rediscover vulnerability-triggering inputs. The grammar evolved through three versions (v3.0 to v3.2), growing from 449 to 475 production rules, and targets the six CVEs identified above across four SQLite releases.

The second objective is to integrate a multi-armed bandit algorithm, specifically the EXP3 algorithm \cite{exp3}, into the grammar rule selection mechanism of the Nautilus fuzzer \cite{nautilus}. EXP3 (Exponential-weight algorithm for Exploration and Exploitation) is designed for adversarial bandit settings where reward distributions may be non-stationary — a natural fit for fuzzing, where the coverage landscape changes as the corpus evolves. The integration requires implementing O(1) weighted sampling via the alias method \cite{loadeddice}, a coverage reward signal derived from new bitmap edges discovered per generated input, and an exponential moving average normalisation scheme to prevent weight compounding over long campaigns.

The third objective is to conduct controlled empirical experiments comparing the bandit-guided and uniform sampling policies across multiple SQLite versions with confirmed CVEs, using a rigorous statistical methodology. The experimental scope encompasses 53 fuzzing campaigns of 30 minutes each across the four target SQLite versions (3.30.1, 3.31.1, 3.32.0, and 3.32.2) and two grammar variants, with crash deduplication by stack hash and statistical comparison using the Mann-Whitney U test \cite{mannwhitney} to assess whether observed differences in crash class counts are statistically significant.

\section{Contributions}

This thesis makes three primary contributions to the area of grammar-based fuzzing for database vulnerability detection. The first contribution is a structural primitives grammar design methodology for SQL-based fuzzing. Rather than encoding specific SQL statements known to trigger bugs, the methodology decomposes SQL into layered non-terminals — \texttt{Schema-Setup}, \texttt{Stress-Query}, \texttt{Validation-Op}, and \texttt{Boundary-Func} — that capture the structural patterns associated with vulnerability classes without contaminating the fuzzer's search with proof-of-concept bias. This compositional design allows the fuzzer to independently reconstruct vulnerability-triggering inputs from primitive combinations, generalise across SQLite versions, and avoid the common failure mode where a grammar is so specific to known bugs that it fails to discover new ones. The grammar's three-version evolution (v3.0 to v3.2, 449 to 475 rules) demonstrates how iterative refinement guided by coverage and crash evidence can progressively improve a structural grammar.

The second contribution is the integration of EXP3 multi-armed bandit adaptive sampling into grammar-based greybox fuzzing, with O(1) weighted sampling via the alias method \cite{loadeddice}. The implementation includes a coverage-derived reward signal, exponential moving average normalisation, and a corrected delta computation that properly credits the bandit for newly discovered bitmap edges rather than compounding weight updates. This constitutes the first integration of EXP3 adaptive sampling into a Nautilus-class grammar fuzzer with a formal evaluation against a controlled baseline.

The third contribution is an empirical finding that challenges the intuitive appeal of adaptive sampling for vulnerability discovery. Across 53 campaigns, uniform random sampling discovers approximately twice as many unique crash classes as bandit-guided sampling: a mean of 8.6 crash classes versus 4.2 for the bandit (p\,≈\,0.01, Mann-Whitney U test \cite{mannwhitney}). This result demonstrates that structural diversity — the breadth of SQL constructs exercised — is more important than coverage exploitation for discovering CVE-class vulnerabilities in SQLite, and provides empirical grounding for the claim that exploration outperforms exploitation in this setting. The result also reveals a previously undocumented failure mode of bandit-guided fuzzing: weight concentration on a small subset of highly rewarded rules collapses input diversity and causes the fuzzer to miss entire vulnerability classes.

\section{Thesis Organisation}

The remainder of this thesis is organised as follows. Chapter~\ref{chap:background} surveys the background literature in four areas: the foundations of coverage-guided greybox fuzzing and grammar-based input generation, the multi-armed bandit formulation and the EXP3 algorithm, the dynamic analysis tools (AddressSanitizer and UndefinedBehaviorSanitizer) used as vulnerability oracles, and the related work on grammar-aware fuzzers and database testing tools most directly relevant to this thesis. Chapter~\ref{chap:method} presents the proposed system: the overall architecture of the grammar-based fuzzer targeting SQLite, the structural primitives grammar design methodology, the integration of EXP3 bandit sampling into the rule selection mechanism, the SQLite harness design and sanitiser oracle configuration, and the crash triage and deduplication pipeline. Chapter~\ref{chap:experiments} reports the experimental evaluation, structured around four research questions covering crash discovery rate, crash class diversity, coverage growth trajectory, and the statistical significance of the bandit-versus-uniform comparison; results are presented for all four SQLite target versions and discussed in light of the original hypotheses. The Conclusion synthesises the key findings, identifies the limitations of the current work, and outlines directions for future research including deeper reinforcement learning integration, grammar generalisation across database engines, and longer-duration campaign studies.
